% Generated. Do not edit manually. Edit resume_content.json instead.
\renewcommand{\ResumeSummary}{
Ph.D. mechanical engineer developing physics-based simulation, digital twins, and engineering automation. Built coupled plant models for NASA habitat research, an experimentally calibrated structural digital twin, reduced-order models, and a neural-network damage surrogate. Professional work includes real-time C/C++ HVAC modeling, SiL integration, and Python/CadQuery automation linking engineering inputs to parametric geometry.
}

\renewcommand{\ResumeSkills}{
\textbf{Digital Twins and Modeling:} Physics-based modeling, Multiphysics coupling, System-of-systems integration, Model calibration, Modal reduction, Guyan condensation\par
\textbf{Data-Driven Simulation:} Neural-network surrogate modeling, Simulation data analysis, Design-space exploration, High-performance computing\par
\textbf{Numerical Software:} Python, C/C++, CadQuery, Numerical algorithms, Custom FEM/XFEM solvers, Simulation workflow automation\par
\textbf{Integration and Validation:} SiL/HiL plant models, Virtual CAN/LIN, Jenkins CI/CD, Bash, PowerShell, Experimental modal analysis\par
}

\renewcommand{\RheemBullets}{
\item \textbf{Parametric Fan-Shroud Generation}
\begin{itemize}[leftmargin=1.3em,topsep=1pt,itemsep=1pt]
\item Developed a Python/CadQuery-based tool that converts fan diameter, coil geometry, component locations, and design clearances into 3D shroud geometry and STEP files, reducing repetitive CAD modeling and accelerating design iteration across product configurations.
\end{itemize}
\item \textbf{Engineering Workflow Automation}
\begin{itemize}[leftmargin=1.3em,topsep=1pt,itemsep=1pt]
\item Built workflows linking parametric design inputs, numerical evaluation, optimization, and CAD output to enable programmatic exploration of engineering alternatives.
\end{itemize}
\item \textbf{Simulation-Driven Design Optimization}
\begin{itemize}[leftmargin=1.3em,topsep=1pt,itemsep=1pt]
\item Developing a framework that uses batch CFD simulations to build a performance dataset across fan-shroud geometry parameters. Applying statistical analysis and predictive modeling to identify promising configurations, establish data-driven design guidelines, and reduce the need for computationally expensive CFD runs during future product development.
\end{itemize}
\item \textbf{Refrigerant-Tube Routing Optimization}
\begin{itemize}[leftmargin=1.3em,topsep=1pt,itemsep=1pt]
\item Developing an FEA-based tool to evaluate natural frequencies and optimize compressor-to-outdoor-coil tubing routes to reduce vibration susceptibility. Incorporating experimental modal testing and model updating to calibrate FEA predictions against measured vibration characteristics.
\end{itemize}
\item \textbf{Motor-Mount Structural Optimization}
\begin{itemize}[leftmargin=1.3em,topsep=1pt,itemsep=1pt]
\item Redesigned a motor-mount structure using FEA, reducing material usage by more than 50\% while maintaining stiffness and vibration requirements. Simplified injection-molded components to reduce secondary manufacturing operations and assembly complexity.
\end{itemize}
}

\renewcommand{\LTTSBullets}{
\item \textbf{Real-Time C/C++ HVAC Plant Model for SiL and HiL}
\begin{itemize}[leftmargin=1.3em,topsep=1pt,itemsep=1pt]
\item Developed a real-time C/C++ executable HVAC plant model for Software-in-the-Loop (SiL) and Hardware-in-the-Loop (HiL) applications, combining finite element solution of transient heat transfer and velocity-pressure-coupled Navier-Stokes equations with lumped-capacitance heater/cooler models, convective thermal-fluid coupling, and a proportional-derivative (PD) HVAC controller model.
\end{itemize}
\item \textbf{C\#/.NET Database and Desktop Application Integration}
\begin{itemize}[leftmargin=1.3em,topsep=1pt,itemsep=1pt]
\item Implemented database read/write operations using SQL and Entity Framework within C\#/.NET Windows Forms (WinForms) applications.
\end{itemize}
\item \textbf{Integrated SiL Simulation Infrastructure and Vehicle Network Virtualization}
\begin{itemize}[leftmargin=1.3em,topsep=1pt,itemsep=1pt]
\item Developed and integrated a Software-in-the-Loop (SiL) infrastructure connecting the simulator, display, virtual CAN/LIN communication, and ECU-testing components to support construction and forestry vehicle software testing.
\end{itemize}
\item \textbf{Jenkins CI/CD and Automated Debugging}
\begin{itemize}[leftmargin=1.3em,topsep=1pt,itemsep=1pt]
\item Developed Jenkins CI/CD pipelines with Bash and PowerShell scripting and automated screen-capture capability, streamlining debugging and test workflows for construction and forestry vehicle software.
\end{itemize}
\item \textbf{Modular GUI for Aerodynamic Data Analysis}
\begin{itemize}[leftmargin=1.3em,topsep=1pt,itemsep=1pt]
\item Developed a modular GUI for aerodynamic data analysis with well-defined input/output interfaces, improving maintainability and reducing computational and memory overhead.
\end{itemize}
\item \textbf{Probability-Based File-Type Detection Algorithm}
\begin{itemize}[leftmargin=1.3em,topsep=1pt,itemsep=1pt]
\item Developed a probability-based file-type detection algorithm for John Deere software, replacing a rigid if-else structure with a scalable, data-driven approach intended to improve performance, reliability, and extensibility.
\end{itemize}
}

\renewcommand{\UTSABullets}{
\item \textbf{XFEM-Informed Neural Network Damage Surrogate}
\begin{itemize}[leftmargin=1.3em,topsep=1pt,itemsep=1pt]
\item Developed a neural-network surrogate trained on high-fidelity XFEM simulations to approximate localized damage behavior in larger structural models.
\end{itemize}
\item \textbf{NASA Modular-Coupled Virtual Testbed (MCVT)}
\begin{itemize}[leftmargin=1.3em,topsep=1pt,itemsep=1pt]
\item Developed and integrated structural, thermal, and fluid plant models for NASA's Modular-Coupled Virtual Testbed (MCVT), a system-of-systems simulation environment for virtual testing of resilient extraterrestrial habitats.
\end{itemize}
\item \textbf{Experimental Structural Digital Twin and Model Updating}
\begin{itemize}[leftmargin=1.3em,topsep=1pt,itemsep=1pt]
\item Created a structural digital twin of an experimental habitat by acquiring impact-hammer and accelerometer data, identifying natural frequencies and mode shapes, using Crypton optical geometry measurements, and calibrating a parameterized beam-element model to match identified modal properties and measured acceleration histories.
\end{itemize}
\item \textbf{Reduced-Order Modeling and High-Order Elements}
\begin{itemize}[leftmargin=1.3em,topsep=1pt,itemsep=1pt]
\item Developed modal-coordinate reduced-order models and applied Guyan static condensation; implemented very high-order finite element formulations for acoustic wave propagation associated with micrometeoroid impacts.
\end{itemize}
\item \textbf{Efficient Scientific Software and Numerical Algorithms}
\begin{itemize}[leftmargin=1.3em,topsep=1pt,itemsep=1pt]
\item Developed computational tools emphasizing reduced arithmetic and memory complexity to support novel numerical formulations and simulation workflows; contributed to a publication on algorithm development.
\end{itemize}
\item \textbf{Custom 3D Multiphysics Finite Element Solvers}
\begin{itemize}[leftmargin=1.3em,topsep=1pt,itemsep=1pt]
\item Developed custom three-dimensional finite element solvers from first principles for structural mechanics, transient heat transfer, and fluid dynamics, including coupled structural-thermal habitat models and interior thermal-fluid simulations.
\end{itemize}
\item \textbf{Thermal Protective Layer Topology Optimization}
\begin{itemize}[leftmargin=1.3em,topsep=1pt,itemsep=1pt]
\item Investigated thermo-mechanical topology optimization of lunar habitat thermal protective layers under different environmental conditions and lunar latitudes.
\end{itemize}
\item \textbf{XFEM, Nonlinear Materials, and Damage Modeling}
\begin{itemize}[leftmargin=1.3em,topsep=1pt,itemsep=1pt]
\item Developed XFEM methods to represent stress fields around three-dimensional voids and implemented nonlinear material behavior incorporating plasticity, strain-rate effects, thermal expansion, and von Mises and Drucker-Prager yield criteria.
\end{itemize}
}

\renewcommand{\OtherExperience}{

}

\renewcommand{\Demonstrations}{
MCVT structural-mechanical model: \url{https://youtu.be/TAA66P_UXIM}\par
Finite element formulation software tool: \url{https://www.linkedin.com/feed/update/urn:li:activity:7122914707790843904/}\par
Three-Minute Thesis: Computational Model for Habitat Structures Beyond Earth: \url{https://rrpress.utsa.edu/items/4d1a04f4-08fd-423d-957f-886608ee12d0}\par
}

\let\ResumeDemonstrations\Demonstrations\renewcommand{\Demonstrations}{}\renewcommand{\OtherExperience}{}
